\begin{table}[H]
\centering
\small
\caption{Что передавалось в модели отбора CatBoost и LogReg}
\label{tab:selector-training-data}
\begin{tabular}{p{0.24\textwidth}p{0.66\textwidth}}
\toprule
\textbf{Компонент} & \textbf{Описание} \\
\midrule
Наблюдение & Одна строка соответствует одному сформированному эпизоду, то есть одному payload после группировки событий. В основном эксперименте используется стратегия \texttt{related\_stop\_cap}. \\
Целевая переменная & \texttt{soft\_label\_3}: 0 --- не комментировать, 1 --- кратко прокомментировать, 2 --- обязательно прокомментировать. Метка получается из эвристической пятибалльной шкалы \texttt{soft\_label\_5}. \\
Числовые признаки & Длина и длительность эпизода, временные разрывы, счётчики классов событий, признаки стандартов, ударов, потерь, продвижения мяча, входа в чужую половину и штрафную, переводов фланга, негативных исходов передач и пространственных переходов. \\
Категориальные признаки & Тип первого и последнего события, тайм, стадия турнира, начальные и конечные относительные зоны, доминирующий переход зоны и другие дискретные описания структуры эпизода. \\
Что не передавалось & Сырые JSON целиком, текст комментария, ответ LLM, идентификаторы событий как самостоятельные признаки и целевые поля. В честных режимах также исключались прямые компоненты формулы score. \\
Разбиение & Train/validation/test формируются по \texttt{match\_id}, чтобы события одного матча не попадали одновременно в обучение и тест. Это снижает риск утечки между соседними эпизодами одного матча. \\
Предобработка & Для LogReg числовые признаки масштабируются, категориальные кодируются. CatBoost получает числовые признаки и категориальные столбцы напрямую через механизм \texttt{cat\_features}. \\
\bottomrule
\end{tabular}
\end{table}
